\begin{thebibliography}{99}
\bibitem{Cautis.Kamnitzer.Morrison} Sabin Cautis, Joel Kamnitzer, and Scott Morrison, Webs and quantum skew Howe duality, Math. Ann. 360 (2014), no. 1-2, 351–390. \url{https://doi.org/10.1007/s00208-013-0984-4}
\bibitem{Elias} Ben Elias, Light ladders and clasp conjectures, 2015, \url{https://arxiv.org/abs/1510.06840}
\bibitem{Fomin.Pylyavskyy} Sergey Fomin and Pavlo Pylyavskyy, Tensor diagrams and cluster algebras, Adv. Math. 300 (2016), 717–787. \url{https://doi.org/10.1016/j.aim.2016.03.030}
\bibitem{Fontaine} Bruce Fontaine, Generating basis webs for SLn, Adv. Math. 229 (2012), no. 5, 2792–2817. \url{https://doi.org/10.1016/j.aim.2012.01.016}
\bibitem{Fraser:2column} Chris Fraser, Webs and canonical bases in degree two, Comb. Theory 3 (2023), no. 3, article no. 11 (26 pages). 
\bibitem{Fraser.Lam.Le} Chris Fraser, Thomas Lam, and Ian Le, From dimers to webs, Trans. Amer. Math. Soc. 371 (2019), no. 9, 6087–6124. \url{https://doi.org/10.1090/tran/7641}
\bibitem{Fraser.Pylyavskyy} Chris Fraser and Pavlo Pylyavskyy, Tensor diagrams and cluster combinatorics at punctures, Adv. Math. 412 (2023), article no. 108796 (83 pages). \url{https://doi.org/10.1016/j.aim.2022.108796}
\bibitem{Fulton:YoungTableaux} William Fulton, Young tableaux. With applications to representation theory and geometry, London Mathematical Society Student Texts, vol. 35, Cambridge University Press, Cambridge, 1997. 
\bibitem{Fulton.Harris} William Fulton and Joe Harris, Representation theory. A first course, Readings in Mathematics, Graduate Texts in Mathematics, vol. 129, Springer-Verlag, New York, 1991. 
\bibitem{Gaetz.Pechenik.Pfannerer.Striker.Swanson:SL4} Christian Gaetz, Oliver Pechenik, Stephan Pfannerer, Jessica Striker, and Joshua P. Swanson, Rotation-invariant web bases from hourglass plabic graphs, 2023,  12501. \url{https://arxiv.org/abs/2306.12501}
\bibitem{Gaetz.Pechenik.Pfannerer.Striker.Swanson:2column} Christian Gaetz, Oliver Pechenik, Stephan Pfannerer, Jessica Striker, and Joshua P. Swanson, Web bases in degree two from hourglass plabic graphs, 2024,  13978. \url{https://arxiv.org/abs/2402.13978}
\bibitem{Hwang.Jang.Oh} Byung-Hak Hwang, Jihyeug Jang, and Jaeseong Oh, A combinatorial model for the transition matrix between the Specht and SL2-web bases, Forum Math. Sigma 11 (2023), article no. e82 (17 pages). \url{https://doi.org/10.1017/fms.2023.79}
\bibitem{Im.Zhu} Mee Seong Im and Jieru Zhu, Transitioning Between Tableaux and Spider Bases for Specht Modules, Algebr. Represent. Theory 25 (2022), no. 2, 387–399. \url{https://doi.org/10.1007/s10468-020-10026-6}
\bibitem{Kim:embedding} Jesse Kim, An embedding of the skein action on set partitions into the skein action on matchings, Electron. J. Combin. 31 (2024), no. 1, article no. 1.17 (17 pages). 
\bibitem{Kim:3row} Jesse Kim, Rotation invariant webs for three row flamingo Specht modules, 2024, \url{https://arxiv.org/abs/2402.11994}
\bibitem{Kim.Rhoades} Jesse Kim and Brendon Rhoades, Set Partitions, Fermions, and Skein Relations, Int. Math. Res. Not. IMRN 2023 (2023), no. 11, 9427–9480. 
\bibitem{Kung.Rota} Joseph P. S. Kung and Gian-Carlo Rota, The invariant theory of binary forms, Bull. Amer. Math. Soc. (N.S.) 10 (1984), no. 1, 27–85. \url{https://doi.org/10.1090/S0273-0979-1984-15188-7}
\bibitem{Kuperberg} Greg Kuperberg, Spiders for rank 2 Lie algebras, Comm. Math. Phys. 180 (1996), no. 1, 109– 151. \url{https://doi.org/10.1007/BF02101184}
\bibitem{Lusztig} G. Lusztig, Canonical bases arising from quantized enveloping algebras, J. Amer. Math. Soc. 3 (1990), no. 2, 447–498. \url{https://doi.org/10.1090/S0894-0347-1990-1035415-6}
\bibitem{Patrias.Pechenik:evacuation} Rebecca Patrias and Oliver Pechenik, Tableau evacuation and webs, Proc. Amer. Math Soc. Ser. B 10 (2023), 341–352. \url{https://doi.org/10.1090/bproc/191}
\bibitem{Patrias.Pechenik.Striker} Rebecca Patrias, Oliver Pechenik, and Jessica Striker, A web basis of invariant polynomials from noncrossing partitions, Adv. Math. 408 (2022), article no. 108603 (33 pages). 
\bibitem{Pechenik:CSP} Oliver Pechenik, Cyclic sieving of increasing tableaux and small Schröder paths, J. Combin. Theory Ser. A 125 (2014), 357–378. \url{https://doi.org/10.1016/j.jcta.2014.04.002}
\bibitem{Petersen.Pylyavskyy.Rhoades} T. Kyle Petersen, Pavlo Pylyavskyy, and Brendon Rhoades, Promotion and cyclic sieving via webs, J. Algebraic Combin. 30 (2009), no. 1, 19–41. \url{https://doi.org/10.1007/s10801-008-0150-3}
\bibitem{Purbhoo.Wu} Kevin Purbhoo and Shelley Wu, Folding rotationally symmetric tableaux via webs, Ann. Comb. 28 (2024), no. 1, 93–119. \url{https://doi.org/10.1007/s00026-023-00648-0}
\bibitem{Rhoades:thesis} Brendon Rhoades, Cyclic sieving, promotion, and representation theory, J. Combin. Theory Ser. A 117 (2010), no. 1, 38–76. \url{https://doi.org/10.1016/j.jcta.2009.03.017}
\bibitem{Rhoades:skein} Brendon Rhoades, A skein action of the symmetric group on noncrossing partitions, J. Algebraic Combin. 45 (2017), no. 1, 81–127. \url{https://doi.org/10.1007/s10801-016-0701-y}
\bibitem{Rhoades:polytabloid} Brendon Rhoades, The polytabloid basis expands positively into the web basis, Forum Math. Sigma 7 (2019), article no. e26 (8 pages). 
\bibitem{Russell.Tymoczko:sl3} Heather M. Russell and Julianna Tymoczko, The transition matrix between the Specht and sl3 web bases is unitriangular with respect to shadow containment, Int. Math. Res. Not. IMRN 2022 (2022), no. 5, 3371–3416. \url{https://doi.org/10.1093/imrn/rnaa290}
\bibitem{Russell.Tymoczko:sl2} Heather M. Russell and Julianna S. Tymoczko, The transition matrix between the Specht and web bases is unipotent with additional vanishing entries, Int. Math. Res. Not. IMRN 2019 (2019), no. 5, 1479–1502. \url{https://doi.org/10.1093/imrn/rnx164}
\bibitem{Sagan} Bruce E. Sagan, The symmetric group. Representations, combinatorial algorithms, and symmetric functions, second ed., Graduate Texts in Mathematics, vol. 203, Springer-Verlag, New York, 2001. 
\bibitem{Sturmfels} Bernd Sturmfels, Algorithms in invariant theory, second ed., Texts and Monographs in Symbolic Computation, SpringerWienNewYork, Vienna, 2008. 
\bibitem{Temperley.Lieb} H. N. V. Temperley and E. H. Lieb, Relations between the “percolation” and “colouring” problem and other graph-theoretical problems associated with regular planar lattices: some exact results for the “percolation” problem, Proc. Roy. Soc. London Ser. A 322 (1971), no. 1549, 251–280. 
\bibitem{Westbury} Bruce W. Westbury, Web bases for the general linear groups, J. Algebraic Combin. 35 (2012), no. 1, 93–107. \url{https://doi.org/10.1007/s10801-011-0294-4}
\end{thebibliography}
